\documentclass[11pt]{article}

\usepackage[margin=1.1in]{geometry}
\usepackage{amsmath,amssymb,amsthm}
\usepackage{mathtools}
\usepackage{stmaryrd}
\usepackage[T1]{fontenc}
\usepackage{palatino}
\usepackage{enumitem}
\usepackage{array}
\usepackage{booktabs}
\usepackage{listings}
\usepackage[protrusion=true,expansion=false]{microtype}
\usepackage[hidelinks]{hyperref}
\setlength{\emergencystretch}{2.5em}

\theoremstyle{definition}
\newtheorem{definition}{Definition}
\newtheorem{principle}{Principle}
\newtheorem{construction}{Construction}
\theoremstyle{plain}
\newtheorem{proposition}{Proposition}
\newtheorem{theorem}{Theorem}
\newtheorem{lemma}{Lemma}
\newtheorem{corollary}{Corollary}
\newtheorem{conjecture}{Conjecture}
\newtheorem{question}{Question}
\theoremstyle{remark}
\newtheorem{remark}{Remark}

\newcommand{\V}{\mathcal{V}}
\newcommand{\sumv}{\oplus}
\newcommand{\bigsumv}{\bigoplus}
\newcommand{\tensorv}{\otimes}
\newcommand{\Bool}{\mathbb{B}}
\newcommand{\Rnn}{\mathbb{R}_{\geq 0}}
\newcommand{\Cx}{\mathbb{C}}
\newcommand{\Cand}{\mathrm{Cand}}
\newcommand{\Out}{\mathrm{Out}}
\newcommand{\Proc}{\mathrm{Proc}}
\newcommand{\scong}{\equiv}
\newcommand{\sem}[1]{\llbracket #1 \rrbracket}
\newcommand{\crisp}[1]{\lceil #1 \rceil}
\newcommand{\AC}{\mathrm{AC}}
\newcommand{\DC}{\mathrm{DC}}
\newcommand{\CC}{\mathrm{CC}}
\newcommand{\LLPO}{\mathrm{LLPO}}
\newcommand{\LPO}{\mathrm{LPO}}
\newcommand{\WKL}{\mathrm{WKL}}
\newcommand{\WWKL}{\mathrm{WWKL}}
\newcommand{\lw}{\leq_{\mathrm{W}}}
\newcommand{\ew}{\equiv_{\mathrm{W}}}
\newcommand{\Wstar}{\star}
\newcommand{\Bai}{\mathbb{N}^{\mathbb{N}}}
\newcommand{\Can}{2^{\mathbb{N}}}
\newcommand{\rbisim}{\sim_{r}}

\lstset{
  basicstyle=\ttfamily\small,
  columns=fullflexible,
  keepspaces=true,
  frame=none,
  xleftmargin=1.4em,
  showstringspaces=false
}

\title{\bfseries Choice Principles as Graded Where-Clauses\\[3pt]
\large One programming paradigm at every altitude of the Weihrauch lattice}

\author{L. Gregory Meredith\\
\small F1R3FLY.io / F1R3FLY Industries\\
\small \texttt{lgreg.meredith@gmail.com}}

\date{Research Note --- August 2026}

\begin{document}
\maketitle

\begin{abstract}
Two companion notes leave a seam between them. One \cite{choicesched} proposes a
Curry--Howard correspondence for nondeterminism in which choice principles are the
types and fair schedulers are the terms, and lists ``the scheduler calculus'' as an
open problem: the term language does not exist. The other \cite{gradedwhere} supplies
a term-level paradigm --- the \texttt{where} clause of a for-comprehension, valued not
in Booleans but in a resolution algebra $\V$, evaluated per match --- and shows that
one syntax covers the Boolean, stochastic, quantum and probabilistic-logic regimes.
This note joins them. We argue that a graded clause is not a choice function but a
\emph{quantifier} in the sense of Escard\'o--Oliva, that a resolver is a
\emph{selection function}, and that the monad morphism between them is exactly the
statement that a resolver realises a clause. Attainability of the induced quantifier
is then a sharp dividing line, and we compute where it falls. We extend the graded
sort with a fixed-point operator --- omitted by design in \cite{gradedwhere} --- and
show that without it every graded clause sits at the bottom of the hypercomputation
stalk. With it, three independent engines of ascent become available: recursion
\emph{depth}, candidate \emph{width}, and \emph{reflection}. The first two are
incomparable, which is the syntactic shadow of the incomparability of $\CC$ and weak
K\"onig's lemma in the Weihrauch lattice, and hence of that lattice being a lattice and
not a chain. We then write down a clause schema for each of
$\LLPO$, $\CC$, $\WKL$, $\WWKL$, $\lim$, $\DC$ and closed choice on Baire space, each
inhabiting the same \texttt{where} slot, so that a scheduler's choice type is readable
off its clause. Full $\AC$ is reachable, by exactly two routes: an \emph{aperture}
channel fed by the environment rather than by the term, and \emph{reflection} --- and
we note that Krivine realises dependent choice with a \texttt{quote} instruction which
the rho calculus has as a term former. We show that the complex-valued regime is
\emph{transverse} to the lattice rather than high in it, that fairness disciplines are
clause combinators rather than typing side-conditions, and that the mechanism buying
$\AC$ is the mechanism that threatens subject reduction --- which is the open hinge of
\cite{agency} in syntactic dress. Every number reported is computed by the
accompanying \texttt{choicesim.py}.
\end{abstract}

\section{Introduction}\label{sec:intro}

\subsection{Two notes and the seam between them}

The note \cite{choicesched} proposes that the mathematical theory of choice is a
declarative language for nondeterminism and the theory of fairness and scheduling its
term language: choice principles are the types, schedulers are the terms, running a
schedule is normalisation. Its dictionary is suggestive and its metatheory is
well-posed --- progress is fairness, subject reduction is the statement that scheduling
cannot bootstrap an oracle. But its second open problem, P2, is the term language
itself: \emph{define the scheduler calculus}. Nothing in the note says what a scheduler
looks like as a piece of syntax a programmer writes.

The note \cite{gradedwhere} was written for a different purpose --- to give syntactic
expression to the resolution of nondeterminism so that a programmer can arrange
quantum interference --- and it arrives at a paradigm. The for-comprehension of rholang
1.4 already carries a \texttt{where} clause. Let its condition take values not in
$\Bool$ but in a commutative semiring $\V$, and let it be evaluated \emph{per match},
once for each candidate assignment of data to patterns. Then one syntax covers
Boolean gating, Gillespie propensities, Born amplitudes, Viterbi decoding, tropical
cost-optimisation and PLN-driven inference; the resolver is the choice of $\V$.

The seam is obvious once stated. \cite{gradedwhere} has a syntax for resolution and no
theory of how much power a given resolution costs. \cite{choicesched} has a theory of
how much power a resolution costs and no syntax. The essay both descend from
\cite{agency} says what the join should look like: the choice type $\chi(c)$ of a cut
is a semantic invariant, not in general computable, but ``approximable from above by
\emph{types} \dots\ the term assignment of the companion note is the effective
approximation of the decoration, which is why a factorisation defined semantically can
nonetheless be certified syntactically.'' A graded clause is exactly such an effective
approximation, and this note builds it.

\subsection{The thesis}

\begin{quote}
\emph{One \texttt{where} slot suffices at every altitude. What changes as one climbs
the choice hierarchy is not the programming construct but the shape of the formula in
it: the presence and polarity of a fixed point, the width of the candidate set it
ranges over, and whether it quantifies through the namespace or through an aperture.
A scheduler's choice type is a syntactic property of its clause.}
\end{quote}

We take the Weihrauch lattice \cite{weihrauch} as the spine rather than the coarse
ladder --- finite choice, $\CC$, $\DC$, $\AC$ --- because the interesting formulae
live at degrees the coarse ladder does not name --- $\LLPO$, $\WWKL$, $\lim$ --- and because two of those
degrees turn out to be the ones a rho programmer writes by accident.

\subsection{Contributions}

\begin{enumerate}[label=C\arabic*.,leftmargin=2.6em,itemsep=3pt]
\item \textbf{The seam, correctly identified.} A graded clause denotes a map
$\Cand \to V$: a weighting of a fibre, not a section of it. Clause, choice principle
and scheduler stand in the relation \emph{family} / \emph{assertion that a section
exists} / \emph{the section} (\S\ref{sec:seam}).
\item \textbf{Attainment as the soundness relation.} A clause is a quantifier
($K$-monad), a resolver is a selection function ($J$-monad), and the morphism
$\varepsilon \mapsto \lambda p.\,p(\varepsilon(p))$ is ``this resolver realises this
clause''. We characterise attainability: it is total exactly when $\sumv$ is idempotent
and the winning candidate carries the multiplicative unit, and it fails for every
genuinely graded clause over $(+,\times)$, where the induced quantifier is an
expectation (\S\ref{sec:attain}, Proposition~\ref{prop:attain}, measured).
\item \textbf{Non-attainment is a lift, not a defect.} Where the quantifier is not
attained, the resolver selects into a larger codomain: distributions for $\Rnn$
(verified by Monte Carlo), amplitudes for $\Cx$. The $\Rnn$ row of the algebra table
is the $\WWKL$ row of the lattice (\S\ref{ssec:wwkl}).
\item \textbf{The graded fixed point, and why nothing climbs without it.} The graded
sort of \cite{gradedwhere} is finitary by design. Proposition~\ref{prop:bottom}: every
fixed-point-free clause has choice type $\bot$. The ascent is the reintroduction of
$\mu$ and $\nu$ (\S\ref{sec:fix}).
\item \textbf{Three engines, two of them incomparable.} Depth, width, reflection.
Depth and width are Weihrauch-incomparable, which is the syntactic reason the order on
scheduler types is a lattice and not a chain (\S\ref{sec:engines}).
\item \textbf{The formulae.} A clause schema per degree, all in one \texttt{where}
slot, with the rho term each sits in (\S\ref{sec:formulae}, Table~\ref{tab:master}).
\item \textbf{Full $\AC$, by exactly two routes.} A ceiling theorem for closed terms
over a presented store, and the two ways past it: the aperture and reflection. Krivine
realises $\DC$ with \texttt{quote}; the rho calculus has \texttt{quote} as a term
former (\S\ref{sec:ac}).
\item \textbf{Interference is transverse.} A resolver with a non-empty interference
deficit is not a choice function at all, so $\Cx$ is off the lattice rather than high
on it (\S\ref{sec:transverse}).
\item \textbf{Fairness disciplines are clause combinators.} The bootstrap trap of
\cite{gradedwhere} is a fairness violation manufactured by the value algebra, and its
repair $\psi \sumv \epsilon$ is weak fairness written as a formula
(\S\ref{sec:fairness}).
\item \textbf{The hinge, sharpened.} The mechanism that buys $\AC$ is the mechanism
that breaks subject reduction. A stratification judgement recovers it, and $\AC$ lives
outside the stratified fragment (\S\ref{sec:sr}).
\end{enumerate}

\subsection{What this note does not do}

It does not prove the Weihrauch degrees claimed for the schemas of
\S\ref{sec:formulae}. Each schema is presented with the reduction it is intended to
realise and with the resolver obligation it imposes; turning ``this clause requires
$\WKL$'' into ``the choice type of this clause is $\ew \WKL$'' requires fixing a
representation of the candidate family as a point of a represented space, which we set
up only informally. It does not present the topos-theoretic setting in which the
Diaconescu argument of \S\ref{ssec:diaconescu} would be a theorem rather than an
analogy. And it does not treat the quantale-graded calculus of \cite{fuzzyrho}, which
is the value-algebra point at which the graded fixed point exists by Knaster--Tarski
and which therefore deserves its own treatment.

%%%%%%%%%%%%%%%%%%%%%%%%%%%%%%%%%%%%%%%%%%%%%%%%%%%%%%%%%%%%%%%%%%%%%%%%%%%%%%%
\section{The seam}\label{sec:seam}

\subsection{What a clause denotes}

We recall the object from \cite{gradedwhere}. A \emph{receipt site} $r$ in a term $P$
is a for-comprehension together with the channels its receipts name. A
\emph{candidate} at $r$ is an assignment $\sigma$ of a distinct datum occurrence to
each pattern of each group such that every pattern matches its assigned datum; write
$\Cand(P,r)$ for the set of candidates and $\Out(r,\sigma)$ for the residual term the
firing produces. A \emph{resolution algebra} is a commutative semiring
$\V = (V,\sumv,\tensorv,0,1)$, with conjunction interpreted by $\tensorv$ and
disjunction by $\sumv$. A graded clause $\psi$ at $r$ denotes
\[
  \sem{\psi}_r : \Cand(P,r) \longrightarrow V .
\]
Candidates are indexed by occurrences, not by terms. The classical interpreter sees
nothing of the value except whether it vanishes: a communication is enabled exactly
when its clause is nonzero, so the crisp language is the $\Bool$-valued fragment and
support adequacy holds by construction.

\subsection{What a choice principle asserts}

A choice principle asserts that a family of nonempty sets admits a section. In the
Weihrauch setting \cite{weihrauch} the canonical form is \emph{closed choice}: $C_X$
takes a closed subset $A \subseteq X$, given negatively (by an enumeration of basic
open sets exhausting its complement), together with the promise $A \neq \emptyset$,
and must return a point of $A$. The degree of $C_X$ depends on $X$: $C_{\mathbb{N}}$ is
countable choice; $C_{\Can} \ew \WKL$; $C_{\Bai}$ is the strongest of the family;
$C_2 \ew \LLPO$ is the weakest nontrivial one. Reducibility $\lw$ orders them and
Weihrauch composition $\Wstar$ --- solve one problem, then, using its answer, the next
--- composes them.

\subsection{The category error, and the division of labour}

It is tempting, and wrong, to try to \emph{write} a choice principle as a clause. A
clause weights a fibre. A choice principle asserts a section of one. No map
$\Cand \to V$ is a section of $\Cand$, whatever $V$ is.

What a clause can do is \emph{determine the family whose section is at issue}, and
therefore certify, syntactically, how hard that section is to produce. That is the
division of labour we adopt throughout, and it is the one \cite{agency} asks for.

\begin{center}
\begin{tabular}{@{}ll@{}}
\toprule
\textbf{Object} & \textbf{Role}\\
\midrule
graded clause $\psi$ & the family: which candidates, weighted how\\
choice principle & the assertion that a section exists, uniformly\\
resolver / scheduler & the section\\
attainment & the resolver realises the clause\\
choice type $\chi$ & the Weihrauch degree of the assertion\\
\bottomrule
\end{tabular}
\end{center}

\begin{remark}[Why the certificate can be syntactic though the degree is not computable]
$\chi$ is undecidable at large \cite{agency}. But the schemas of
\S\ref{sec:formulae} bound it from above by inspection: a fixed-point-free clause over
a finite candidate set cannot require more than a definable section, whatever else it
does. Certification from above is all a type discipline ever offers, and it is what
makes ``the oracle it imports is written on its face'' \cite{choicesched} an
implementable slogan.
\end{remark}

%%%%%%%%%%%%%%%%%%%%%%%%%%%%%%%%%%%%%%%%%%%%%%%%%%%%%%%%%%%%%%%%%%%%%%%%%%%%%%%
\section{Quantifiers, selections, and attainment}\label{sec:attain}

\subsection{The two monads}

Escard\'o and Oliva \cite{escardo} study two monads on a fixed codomain $R$. The
\emph{quantifier} (continuation) monad
\[
  K_R X = (X \to R) \to R,
\]
and the \emph{selection} monad
\[
  J_R X = (X \to R) \to X .
\]
There is a monad morphism $\overline{(-)} : J_R \to K_R$ given by
\[
  \overline{\varepsilon}(p) \;=\; p(\varepsilon(p)),
\]
and a quantifier $\phi$ is \emph{attainable} when $\phi = \overline{\varepsilon}$ for
some selection function $\varepsilon$. The paradigm example of an attainable
quantifier is $\sup$, attained by $\arg\max$; the paradigm example of a non-attainable
one is expectation.

The identification we propose is immediate once the two monads are on the table.

\begin{definition}[The quantifier of a clause]\label{def:quant}
Let $\psi$ be a graded clause at a receipt site $r$ over $\V$. Its \emph{induced
quantifier} is $\phi_\psi : (\Cand(P,r) \to V) \to V$,
\[
  \phi_\psi(p) \;=\; \bigsumv_{\sigma \in \Cand(P,r)} \sem{\psi}_r(\sigma)
  \tensorv p(\sigma) .
\]
\end{definition}

\noindent
Here $p$ is a \emph{payoff}: an assessment, in $\V$, of what each candidate leads to.
The graded modality supplies the canonical payoff, since
$\sem{\langle K \rangle_\V \psi'}(P) = \sum_{r,\sigma} \sem{\psi_r}(\sigma) \tensorv
\sem{\psi'}(\Out(r,\sigma))$ is exactly $\phi_{\psi_r}$ applied to the payoff
$\sigma \mapsto \sem{\psi'}(\Out(r,\sigma))$. So the graded modality of
\cite{gradedwhere} \emph{is} the clause's quantifier eating the value of the future.

\begin{definition}[A resolver realises a clause]\label{def:realise}
A \emph{resolver} at $r$ is a selection function
$\varepsilon : (\Cand(P,r) \to V) \to \Cand(P,r)$. It \emph{realises} $\psi$ when
$\overline{\varepsilon} = \phi_\psi$, that is, when for every payoff $p$,
\[
  p(\varepsilon(p)) \;=\; \bigsumv_{\sigma} \sem{\psi}_r(\sigma) \tensorv p(\sigma).
\]
\end{definition}

This is the soundness relation the seam needed. It says: the value the clause assigns
to the site as a whole is the value actually obtained by the candidate the resolver
picks. A resolver that realises its clause is not merely consistent with it; it
\emph{achieves} it.

\subsection{When is a clause realisable at all?}

\begin{proposition}[Attainability]\label{prop:attain}
Let $\psi$ be a clause over $\V$ with finite candidate set.
\begin{enumerate}[label=(\roman*),leftmargin=2.2em,itemsep=2pt]
\item If $\sumv$ is idempotent and $\psi$ is crisp --- that is, $\sem{\psi}(\sigma) \in
\{0,1\}$ for every $\sigma$ --- then $\phi_\psi$ is attained by a definable selection
function, namely the $\sumv$-optimal candidate of the support.
\item If $\sumv$ is idempotent and $\psi$ is graded, then $\phi_\psi$ is attained at a
payoff $p$ if and only if the $\sumv$-winning candidate $\sigma^\ast$ carries the
multiplicative unit, $\sem{\psi}(\sigma^\ast) = 1$.
\item If $\sumv$ is not idempotent --- in particular for $\V = \Rnn$ or $\Cx$ --- then
$\phi_\psi$ is an expectation and is not attained by any selection function into
$\Cand$, except in the degenerate case of a singleton support.
\end{enumerate}
\end{proposition}

\begin{proof}[Proof sketch]
(i) With $\sumv$ idempotent, $\phi_\psi(p) = \bigsumv_{\sigma \in \mathrm{supp}(\psi)}
p(\sigma)$, which for $\sumv = \max$, $\min$ or $\vee$ is a value of $p$ on the
support; take $\varepsilon$ to return the candidate achieving it.
(ii) $\phi_\psi(p) = \bigsumv_\sigma \sem{\psi}(\sigma)\tensorv p(\sigma)$. If the
winner carries $1$ this equals $p(\sigma^\ast)$; if it carries $v \neq 1$ the value is
$v \tensorv p(\sigma^\ast)$, which is not in general in the image of $p$.
(iii) $\sum_\sigma \sem{\psi}(\sigma) p(\sigma)$ is a strict convex-type combination
whenever two candidates carry nonzero weight, hence lies strictly between two values of
$p$ and is not among them, generically.
\end{proof}

The proposition is sharp enough to be measured, and we measured it. Over $20{,}000$
random payoffs per row (\texttt{choicesim.py}, E1):

\begin{center}
\begin{tabular}{@{}lr@{}}
\toprule
\textbf{Algebra and clause} & \textbf{attained}\\
\midrule
$\Bool$, crisp $\psi$ & $100.00\%$\\
Viterbi $(\max,\times)$, crisp $\psi$ & $100.00\%$\\
Viterbi $(\max,\times)$, graded $\psi$ & $73.83\%$\\
tropical $(\min,+)$, crisp $\psi$ & $100.00\%$\\
tropical $(\min,+)$, graded $\psi$ & $0.00\%$\\
$\Rnn$, crisp $\psi$ & $0.00\%$\\
$\Rnn$, graded $\psi$ & $0.00\%$\\
\bottomrule
\end{tabular}
\end{center}

\noindent
and the intermediate Viterbi figure is not noise: over the same payoffs, ``$\phi_\psi$
is attained'' agrees with ``the winning candidate carries the unit'' on $100.00\%$ of
cases, which is clause (ii) verified.

\begin{remark}[Restriction versus discount]\label{rem:restrict}
Clause (ii) has a reading worth keeping. Over an idempotent $\sumv$, a clause that
merely \emph{restricts} --- rules candidates in and out --- is realisable by a pure
term; a clause that \emph{discounts} --- prefers one admissible candidate to another
--- is not. Gating is free. Preference costs.
\end{remark}

\subsection{Non-attainment is a lift}\label{ssec:lift}

Clause (iii) looks like bad news for the whole enterprise, since $\Rnn$ and $\Cx$ are
the two rows anyone cares about. It is not. The standard response to a non-attainable
quantifier is to enlarge the codomain of the selection function, and the enlargements
are exactly the resolvers already in use.

\begin{itemize}[leftmargin=1.4em,itemsep=2pt]
\item $\V = \Rnn$. Lift $\varepsilon$ to select into $\mathcal{D}(\Cand)$,
distributions over candidates, taking $p \mapsto$ the normalised clause. Then
$\mathbb{E}[p(\varepsilon(p))] = \phi_\psi(p) / \phi_\psi(\mathbf{1})$, so attainment
holds \emph{in expectation}. Verified by Monte Carlo over $400{,}000$ draws: measured
$0.590826$ against target $0.590000$, and $0.390489$ against $0.390000$
(\texttt{choicesim.py}, E1).
\item $\V = \Cx$. No lift into distributions works, because the coefficients are not
positive. The Born rule is a lift into $\mathcal{D}(\Cand)$ \emph{after} a quadratic
map, and \S\ref{sec:transverse} shows why that is a different kind of object.
\end{itemize}

\begin{principle}[The lift is the oracle]\label{prin:lift}
The extra structure a resolver must have, beyond being a function into the candidate
set, is exactly the failure of attainment of its clause's quantifier. Where the clause
is crisp over an idempotent $\sumv$, the scheduler is a pure term and imports nothing.
Where it discounts, the scheduler must optimise. Where $\sumv$ is additive, the
scheduler must randomise. This is Principle~\ref{prin:complementary}'s first payoff:
the altitude is legible in the algebra before any fixed point is written.
\end{principle}

\subsection{The product is the composition}\label{ssec:product}

Escard\'o--Oliva's binary product of selection functions,
\[
  (\varepsilon_0 \tensorv \varepsilon_1)(p) = (a_0, a_1), \quad
  a_0 = \varepsilon_0\big(\lambda x.\, p(x, \varepsilon_1(\lambda y.\,p(x,y)))\big),
  \quad a_1 = \varepsilon_1(\lambda y.\, p(a_0,y)),
\]
is their characteristic operation, and its infinite iteration is bar recursion. Three
things now want to be the same operation: this product; Weihrauch composition
$\Wstar$, ``solve one, then, using its answer, the next'' \cite{weihrauch}, which
\cite{choicesched} nominates as the cut of the logic of choice; and the sequential
chaining of two graded receipt sites through $\langle K \rangle_\V$, in which the
second site's clause is evaluated against the term the first left behind.

\begin{conjecture}[One operation, three names]\label{conj:one}
For clauses whose quantifiers are attainable, the Escard\'o--Oliva product of the
realising resolvers, the Weihrauch composition of the choice types, and the chaining of
the graded sites through $\langle K \rangle_\V$ agree. Consequently $\DC$ is the
$\omega$-fold product, bar recursion is its realiser, and the graded modality is the
syntax of both.
\end{conjecture}

We verified the arithmetic half over Viterbi: the product and the chained sites select
the same pair on $20{,}000$ random payoff tables, with $0$ disagreements. That
agreement is a consistency check rather than a discovery --- but it is not vacuous, and
we made sure it is not. The identification requires the chaining to aggregate in the
\emph{same} algebra the selection function optimises in. Chaining with $\sumv = +$
while still selecting by $\arg\max$ --- backward induction over expectation rather than
over the maximum --- disagrees on $4{,}483$ of the same $20{,}000$ tables
(\texttt{choicesim.py}, E2). So the identification has content, and the content is that
the algebra of the clause and the discipline of the scheduler must match.

%%%%%%%%%%%%%%%%%%%%%%%%%%%%%%%%%%%%%%%%%%%%%%%%%%%%%%%%%%%%%%%%%%%%%%%%%%%%%%%
\section{The graded fixed point}\label{sec:fix}

\subsection{Why the grammar as it stands cannot climb}

The graded sort of \cite{gradedwhere} is
\[
  \psi ::= \crisp{\beta} \mid v \mid \psi \tensorv \psi \mid \psi \sumv \psi
  \mid \langle K \rangle_\V \psi \qquad (v \in V),
\]
with the crisp sort $\beta$ carrying full Boolean structure, the spatial connectives,
the modality and fixed points. The graded sort deliberately has neither negation nor a
fixed point. The exclusion was principled --- $\Cx$ has no order, so Knaster--Tarski is
unavailable and a graded $\nu$ is a convergence question rather than a definition ---
but it has a consequence the note did not draw.

\begin{proposition}[Everything finitary sits at the bottom]\label{prop:bottom}
Let $\psi$ be a graded clause containing no fixed point, at a receipt site whose
candidate set is finite. Then $\phi_\psi$ is computable from the store, and the
resolver realising it (under Proposition~\ref{prop:attain}(i), or its lift under
\S\ref{ssec:lift}) imports no oracle: $\chi = \bot$.
\end{proposition}

\begin{proof}[Proof sketch]
$\crisp{\beta}$ is decidable against a finite term by hypothesis on the crisp fragment;
scalars are constants; $\tensorv$ and $\sumv$ are pointwise; and
$\langle K \rangle_\V$ applied a syntactically bounded number of times ranges over a
finite set of successors, since each site has finitely many candidates and the nesting
depth of the modality in $\psi$ is finite. So $\sem{\psi}_r$ is a computable finite
vector, and choice from a finite presented family is definable outright.
\end{proof}

\begin{corollary}
The whole of \cite{gradedwhere} --- including its Shor construction, whose extractable
fragment is characterised there as the modality-free clauses with isometric values ---
lives at the base of the hypercomputation stalk. The extractable fragment and the
constructively-choosable fragment coincide.
\end{corollary}

That coincidence is worth pausing on. A clause is compilable to a circuit exactly when
its resolution is definable; the reason the ZX pipeline of \cite{cwrhozx} can consume
the fragment it consumes is that the fragment demands no choice. Extraction and
constructivity are the same constraint seen from two ends.

\subsection{The grammar, extended}

\begin{definition}[Graded formulae with fixed points]\label{def:grammar}
Fix a resolution algebra $\V$ and a set of formula variables $X$. The graded sort is
\[
  \psi ::= \crisp{\beta} \mid v \mid X \mid \psi \tensorv \psi \mid \psi \sumv \psi
  \mid \langle K \rangle_\V \psi \mid \langle @ \rangle_\V \psi
  \mid \mu X.\,\psi \mid \nu X.\,\psi ,
\]
subject to the \emph{guardedness} condition that every free occurrence of $X$ in
$\mu X.\psi$ or $\nu X.\psi$ lies under a $\langle K \rangle_\V$ or a
$\langle @ \rangle_\V$. The new modality $\langle @ \rangle_\V \psi$ is the graded
namespace connective: evaluated at a candidate binding a name $y$, it returns the value
of $\psi$ at the process $*y$ quoted by that name.
\end{definition}

Guardedness is what makes the unfolding a genuine iteration of the transfer operator
$T$ of \cite{gradedwhere}, rather than a self-reference with no dynamics. The two
polarities have distinct readings.

\begin{itemize}[leftmargin=1.4em,itemsep=3pt]
\item $\mu X.\psi$ is a \emph{search}: the least fixed point, the value accumulated by
those unfoldings that terminate. Operationally, ``keep unfolding until the body
contributes something nonzero''. A $\mu$ whose unfolding depth is not bounded by the
term is an unbounded search, and unbounded search is where $\CC$ enters.
\item $\nu X.\psi$ is a \emph{survival} condition: the greatest fixed point, the value
carried by unfoldings that never fail. Operationally, ``this holds along some
infinite path''. A $\nu$ over an infinitely deep candidate tree is a statement about
branches, and statements about branches of infinite trees are where K\"onig's lemma
enters.
\end{itemize}

\subsection{When the fixed point exists}

This is where the value algebra stops being a free parameter.

\begin{proposition}[Existence, by algebra]\label{prop:exists}
Let $\Psi$ be the guarded functional on $\V$-valued assignments determined by
$\psi$.
\begin{enumerate}[label=(\roman*),leftmargin=2.2em,itemsep=2pt]
\item If $\V$ carries a complete order compatible with $\sumv,\tensorv$ --- a quantale,
as in \cite{fuzzyrho} --- then $\mu$ and $\nu$ exist by Knaster--Tarski, and no budget
is required.
\item If $\V = \Rnn$, $\Psi$ is monotone but the lattice is not complete;
$\mu X.\psi$ exists as a supremum in $[0,\infty]$ and $\nu X.\psi$ requires a
convergence hypothesis. Sufficient: the spectral radius of the transfer operator
restricted to the reachable set is $< 1$.
\item If $\V = \Cx$ there is no order at all. Neither fixed point is defined by
approximation from below or above; both must be given as solutions of a linear
system, and existence is a spectral condition on $T$.
\end{enumerate}
\end{proposition}

\begin{remark}[The budget is the honest default]\label{rem:budget}
Rholang 1.4 already reins in infinitary property checks with a token budget
\cite{omnibus}. The budget is precisely the retreat from (i) to
Proposition~\ref{prop:bottom}: a budgeted $\mu$ or $\nu$ is a bounded unfolding, hence
finitary, hence at $\bot$. This gives the implementation a dial that is also a
\emph{semantic} dial --- spend budget, gain altitude --- and it means a shard can
refuse to run a clause whose declared altitude exceeds its policy simply by capping
unfoldings.
\end{remark}

\begin{principle}[The two gradings are not orthogonal]\label{prin:complementary}
There are two gradings in play: the strength of the choice principle (a Weihrauch
degree, an altitude in the stalk of \cite{agency}) and the value algebra of the clause.
They are transverse in general --- $\Cx$ is not ``stronger choice'', it is a different
resolver --- but the fixed point couples them, because whether a fixed point exists at
all is a property of the algebra's order structure
(Proposition~\ref{prop:exists}). The altitude a clause certifies is a function of
formula shape \emph{and} algebra.
\end{principle}

%%%%%%%%%%%%%%%%%%%%%%%%%%%%%%%%%%%%%%%%%%%%%%%%%%%%%%%%%%%%%%%%%%%%%%%%%%%%%%%
\section{Three engines of ascent}\label{sec:engines}

With the fixed point available, exactly three things can make a clause's family
non-presented. We name them and then show that two of them are incomparable.

\begin{definition}[The engines]\label{def:engines}
Let $\psi$ be a graded clause at a site $r$ of a term $P$.
\begin{itemize}[leftmargin=1.4em,itemsep=3pt]
\item \emph{Depth.} The \emph{unfolding rank} of $\psi$ is the least $n$ such that
$\psi$'s fixed points stabilise within $n$ applications of the transfer operator,
or $\omega$ if no such $n$ exists. Depth $\omega$ with bounded branching is the
\emph{depth engine}.
\item \emph{Width.} The \emph{candidate width} of $r$ is $\sup |\Cand(P',r)|$ over
terms $P'$ reachable from $P$. Width $\omega$ at bounded depth is the \emph{width
engine}.
\item \emph{Reflection.} $\psi$ uses the \emph{reflection engine} when it contains
$\langle @ \rangle_\V$ applied to a name bound by the site's own receipts, so that the
family it indexes is determined by processes received at run time rather than by
sub-terms of $P$.
\end{itemize}
\end{definition}

Each engine has a plain operational meaning in rholang. Depth $\omega$ is a recursive
process --- built, since core rho has no primitive replication, by the reflection
idiom, a receipt that re-installs itself by dropping its own quotation. Width $\omega$
is a channel on which unboundedly many messages accumulate: a producer under recursion
whose consumer is slower, or a persistent producer. Reflection is a clause that talks
about what it just received.

\begin{proposition}[Depth and width are incomparable]\label{prop:incomparable}
The depth engine at bounded width realises closed choice on a compact space; the width
engine at bounded depth realises closed choice on the naturals. These are
Weihrauch-incomparable: $C_{\mathbb{N}} \not\lw C_{\Can}$ and
$C_{\Can} \not\lw C_{\mathbb{N}}$ \cite{weihrauch}.
\end{proposition}

\begin{proof}[Proof sketch]
Bounded width $k$ and unbounded depth present the candidate family as an infinite,
finitely-branching tree whose nodes are the successive sites' candidates and whose
branches are the runs a $\nu$ quantifies over. A $\nu$-clause asserts the existence of
an infinite branch on which the body never vanishes, and by guardedness the set of such
branches is closed in $k^{\mathbb{N}}$; choosing one is $C_{k^{\mathbb{N}}} \ew \WKL$.
Bounded depth and unbounded width present a single site's candidate set as a subset of
$\mathbb{N}$, given negatively by the clause's vanishing; a $\mu$-clause searches it,
and choosing an element is $C_{\mathbb{N}}$. The incomparability is standard.
\end{proof}

\begin{remark}[Why the order on scheduler types is a lattice]\label{rem:lattice}
\cite{choicesched} insists that the Weihrauch lattice is a lattice and not a chain, and
calls this ``the declarative shadow of the jumble-not-tower ontology'' of
\cite{agency}. Proposition~\ref{prop:incomparable} gives the shadow a syntactic
source: two independent axes of a formula --- how deep it recurses and how wide it
looks --- generate incomparable degrees. A programmer can make a scheduler stronger in
two ways that do not compare, and the non-comparison is visible in the clause without
computing anything.
\end{remark}

%%%%%%%%%%%%%%%%%%%%%%%%%%%%%%%%%%%%%%%%%%%%%%%%%%%%%%%%%%%%%%%%%%%%%%%%%%%%%%%
\section{The formulae}\label{sec:formulae}

We now write a clause schema for each degree. Throughout, the surface form is the same:
\begin{lstlisting}
for( ptrn <- chan where psi ) P
\end{lstlisting}
and only \texttt{psi} changes. We give, for each degree, the schema, the shape of term
it sits in, and the obligation it places on the resolver.

\subsection{The bottom: finite presented choice}\label{ssec:finite}

\[
  \psi \;::=\; \crisp{\beta} \mid v \mid \psi \tensorv \psi \mid \psi \sumv \psi
  \mid \langle K \rangle_\V^{\leq n} \psi
\]
--- no fixed point, or a budgeted one; finite candidate width. By
Proposition~\ref{prop:bottom} the family is presented as a list, and by
Proposition~\ref{prop:attain}(i) a crisp clause over an idempotent $\sumv$ is realised
by a definable selection function. This is the Martin-L\"of point of
\cite{choicesched}: the proof \emph{is} the choice function, and no scheduler
primitive is needed. Every clause in \cite{gradedwhere} is of this shape.

\emph{Resolver obligation:} none. Any deterministic tie-break realises it.

\subsection{$\LLPO \ew C_2$: the negatively-presented binary commit}\label{ssec:llpo}

The weakest nontrivial degree is choice on two points, and it is reachable with two
messages and no fixed point at all --- provided the crisp guards are
\emph{semi-decidable complements} rather than decidable predicates.

\begin{lstlisting}
new a, b in
  ( search(f, a) | search(g, b)
  | for( x <- a ; y <- b  where  [halts(f)] (+) [halts(g)] ) P )
\end{lstlisting}

\noindent
where \texttt{search} produces on its channel exactly if its argument's search
terminates, and the clause's two crisp conjuncts are the corresponding
$\Sigma^0_1$ statements. The promise ``not both'' is the term's invariant; the clause's
support is therefore never empty, and a resolver must commit to one side without being
able to decide which side is safe.

We checked both halves computationally (\texttt{choicesim.py}, E3). Over
$3 \times 40 \times 60$ instance/stage pairs the support was empty $0$ times: the
family is genuinely a family of nonempty sets. And for every stage-$k$ selector,
$k = 0,\dots,7$, under both tie-breaking preferences, the diagonal instance --- which
produces at stage $k+1$ on exactly the side the selector commits to --- defeats it.
So the clause is inhabited but has no stage-uniform computable section.

\emph{Resolver obligation:} $\LLPO$. This is the degree of ``choose which of two
$\Sigma^0_1$ statements fails''.

\begin{remark}[Rho programmers write this by accident]\label{rem:accident}
The term above is an ordinary race between two producers whose production is
conditional on a semi-decidable event. Nothing exotic is involved and nothing in the
current language flags it. That a routine racing idiom sits strictly above the
computable degrees is the clearest argument for wanting the type in the first place.
\end{remark}

\subsection{$C_{\mathbb{N}}$: the width engine}\label{ssec:cn}

\[
  \psi \;=\; \mu X.\, \big( \crisp{\beta} \sumv \langle K \rangle_\V X \big)
\]
at a site whose channel accumulates unboundedly many data. The $\mu$ is an unbounded
search through the candidate set for one satisfying $\beta$, where $\beta$ is again
only semi-decidable, so the set is given negatively: a candidate is excluded when it is
\emph{seen} to fail, and never positively certified in finite time.

\begin{lstlisting}
new c in
  ( enumerate(c)                            // unboundedly many data on c
  | for( x <- c  where  mu X. ([not_excluded(x)] (+) <K>_V X) ) P )
\end{lstlisting}

\emph{Resolver obligation:} $C_{\mathbb{N}}$ --- countable choice, an unbounded search
that terminates because the family is promised nonempty, with no bound computable in
advance on how long it takes.

\subsection{$\WKL \ew C_{\Can}$: the depth engine}\label{ssec:wkl}

\[
  \psi \;=\; \nu X.\, \big( \crisp{\beta} \tensorv \langle K \rangle_\V X \big)
\]
at a site of bounded width --- two candidates suffice --- under a recursive process, so
that the candidate tree is infinite and finitely branching. The $\nu$ asserts the
existence of an infinite run along which $\beta$ never fails: exactly an infinite
branch of an infinite binary tree.

\begin{lstlisting}
new c in
  ( bits(c)                                 // two data at a time, forever
  | for( x <- c  where  nu X. ([ok(x)] (x) <K>_V X) ) P )
\end{lstlisting}

\emph{Resolver obligation:} $\WKL$ --- compactness. The resolver must produce a branch
of a tree all of whose finite paths extend, which is not computable from the tree.

\begin{remark}[The polarity of the fixed point is the polarity of the principle]
$\mu$ searches and $\nu$ survives; $C_{\mathbb{N}}$ searches an unbounded set and
$\WKL$ selects an infinite path. The syntactic distinction between the two fixed
points and the semantic distinction between the two incomparable choice principles are
the same distinction, and Proposition~\ref{prop:incomparable} is that observation made
into a claim.
\end{remark}

\subsection{$\WWKL$: the $\Rnn$ row, and where the two axes touch}\label{ssec:wwkl}

Take the $\WKL$ schema and replace the crisp $\beta$ by a genuinely graded value, with
$\V = \Rnn$ and the clause normalised so that the total over a level is $1$:
\[
  \psi \;=\; \nu X.\, \big( w(\cdot) \tensorv \langle K \rangle_\V X \big),
  \qquad w : \Cand \to \Rnn .
\]
The set of surviving branches is now not merely nonempty but of \emph{positive
measure}, and the resolver's obligation weakens accordingly: it need only find a branch
in a positive-measure closed set. That is $\WWKL$, positive-measure choice, which
Brattka, Gherardi and H\"olzl identify as the degree of Las Vegas computability
\cite{brattkalasvegas} --- randomised computation that never errs but may not
terminate.

This is the one place the two gradings of Principle~\ref{prin:complementary} genuinely
touch. The lift of \S\ref{ssec:lift} --- selecting into distributions over candidates
rather than into candidates, which is forced by Proposition~\ref{prop:attain}(iii) ---
is precisely the randomness that $\WWKL$ measures.

\begin{proposition}[The stochastic row is the randomised degree]\label{prop:wwkl}
A clause over $\Rnn$ with a $\nu$ of unbounded depth, normalised and with positive
total mass at every level, has choice type $\WWKL$, and the resolver realising it is
the distributional lift of \S\ref{ssec:lift}: sampling.
\end{proposition}

\noindent
Read backwards, this says something about the Gillespie regime of \cite{weightedgslt}:
a stochastic simulator is not a weaker classical scheduler but a scheduler at a
specific, nonzero altitude, and the altitude is the price of the random bits.

\subsection{$\lim$: plasticity without a modulus}\label{ssec:lim}

The forager of \cite{gradedwhere} maintains a belief whose value moves as evidence
accumulates: the clause is evaluated against the current term, so its value at a site
depends on what earlier firings did. Suppose the sequence of values converges --- the
belief settles --- but no modulus of convergence is computable from the term. Then a
resolver that must act on the \emph{limit} value is computing a limit of a computable
sequence, and $\lim$ is Weihrauch-equivalent to the Turing jump \cite{weihrauch}.

\[
  \psi \;=\; \nu X.\, \big( \mathrm{pow}(s,c) \tensorv \langle K \rangle_\V X \big)
\]
with $\mathrm{pow}(s,c) = s \cdot c$ the PLN strength-times-confidence value of
\cite{gradedwhere}, $c = n/(n+k)$, and $n$ the evidence count accumulated by firing.

\emph{Resolver obligation:} $\lim$, if the resolver must act on the limit;
$\bot$ if it acts on the current approximation. The distinction is the whole content:
\emph{learning is free, but knowing that you have finished learning is the jump.}

\begin{remark}[The personality parameter, revisited]
\cite{gradedwhere} finds the PLN scepticism parameter $k$ has an interior metabolic
optimum. In the present reading, $k$ is a \emph{modulus policy}: it trades the rate at
which confidence accumulates against the risk of acting on an unsettled value. That an
optimum exists is the statement that neither $\bot$ nor $\lim$ is the right altitude to
run at, and that the environment prices the compromise.
\end{remark}

\subsection{$\DC$: the $\omega$-fold product}\label{ssec:dc}

Dependent choice requires no new syntax whatever. Per-match evaluation against the
current term already makes a site's clause depend on the outcomes of earlier sites,
which is exactly the defining feature of $\DC$: later fibres depend on earlier
selections.

\begin{proposition}[Plasticity is dependent choice]\label{prop:dc}
A graded clause evaluated per match against the running term, at a sequence of sites
chained through $\langle K \rangle_\V$, presents the family
$\prod_{i \in \omega} X_i$ with $X_{i+1}$ depending on the selection made at $X_i$.
Its resolver is the $\omega$-fold Escard\'o--Oliva product of the per-site resolvers,
i.e.\ bar recursion.
\end{proposition}

\noindent
This is Conjecture~\ref{conj:one} instantiated at $\omega$, and it is the cleanest
alignment in the note: \cite{choicesched}'s row ``countable/dependent choice $\DC$ ---
bar recursion; scheduler with memory and lookahead'' is satisfied by a feature
\cite{gradedwhere} already has and describes as a convenience
(``plasticity is free'').

\subsection{$C_{\Bai}$: both engines together}\label{ssec:baire}

Unbounded depth \emph{and} unbounded width --- a $\nu$ over a tree with infinite
branching --- presents the family as a closed subset of Baire space, and
$C_{\Bai}$ is the strongest of the closed-choice family.

\[
  \psi \;=\; \nu X.\, \big( \crisp{\beta} \tensorv \mu Y.(\langle K \rangle_\V
  (X \sumv Y)) \big)
\]
--- the inner $\mu$ searching an unbounded width at each level, the outer $\nu$
surviving unbounded depth.

\emph{Resolver obligation:} $C_{\Bai}$. This is the ceiling for closed terms, as
\S\ref{ssec:ceiling} makes precise.

\subsection{The master table}\label{ssec:table}

\begin{table}[t]
\centering
\small
\begin{tabular}{@{}lllll@{}}
\toprule
\textbf{Degree} & \textbf{Clause shape} & \textbf{Depth} & \textbf{Width} &
\textbf{Resolver must}\\
\midrule
$\bot$ (definable) & no fixed point & $<\omega$ & $<\omega$ & tie-break\\
$\bot$ (optimising) & graded, idempotent $\sumv$ & $<\omega$ & $<\omega$ &
optimise\\
$\LLPO \ew C_2$ & crisp, $\Sigma^0_1$ guards & $<\omega$ & $2$ & commit blind\\
$C_{\mathbb{N}}$ & $\mu X.(\crisp{\beta} \sumv \langle K\rangle_\V X)$ & $<\omega$ &
$\omega$ & search unboundedly\\
$\WKL \ew C_{\Can}$ & $\nu X.(\crisp{\beta} \tensorv \langle K\rangle_\V X)$ &
$\omega$ & $<\omega$ & be compact\\
$\WWKL$ & as $\WKL$, $\V = \Rnn$, normalised & $\omega$ & $<\omega$ & randomise\\
$\lim$ & plastic, acting on the limit & $\omega$ & $<\omega$ & take a limit\\
$\DC$ & chained plastic sites & $\omega$ & any & bar-recurse\\
$C_{\Bai}$ & $\nu$ over $\mu$ & $\omega$ & $\omega$ & choose on Baire space\\
$\AC$ & aperture, or $\langle @ \rangle_\V$ under $\nu$ & --- & --- & be a choice
function\\
\bottomrule
\end{tabular}
\caption{One \texttt{where} slot at every altitude. Only the formula changes.}
\label{tab:master}
\end{table}

%%%%%%%%%%%%%%%%%%%%%%%%%%%%%%%%%%%%%%%%%%%%%%%%%%%%%%%%%%%%%%%%%%%%%%%%%%%%%%%
\section{Reaching full $\AC$}\label{sec:ac}

\subsection{The ceiling for closed terms}\label{ssec:ceiling}

First the obstruction, stated so that the routes past it can be seen to be the only
ones.

\begin{proposition}[Ceiling]\label{prop:ceiling}
Let $P$ be a closed term all of whose channels are fed by sub-terms of $P$, and let
$\psi$ be a graded clause containing no $\langle @ \rangle_\V$ under a fixed point.
Then the union of the candidate sets arising at any site over any run of $P$ is
recursively enumerable from $P$, the family is presented, and
$\chi \lw C_{\Bai}$.
\end{proposition}

\begin{proof}[Proof sketch]
Candidates are assignments of datum occurrences to patterns. Data arise only by
reduction from $P$, and the reduction relation of the rho calculus is r.e.; names are
quoted processes, so the namespace generated by $P$ is r.e. as well. Hence the total
candidate family is an r.e.-indexed family of subsets of an r.e.\ set, given negatively
by the clause's vanishing --- a closed subset of a computably admissible, countably
based space. Closed choice on such a space reduces to $C_{\Bai}$.
\end{proof}

So a closed rho term, however baroque, cannot demand more than choice on Baire space.
The set-theoretic $\AC$ --- ``survey an unindexed contention and select coherently''
\cite{choicesched} --- needs a family that is not presented at all. There are exactly
two ways for a clause to index one.

\subsection{Route E: the aperture}\label{ssec:aperture}

The hypothesis of Proposition~\ref{prop:ceiling} that does most of the work is ``all of
whose channels are fed by sub-terms of $P$''. Drop it.

\begin{definition}[Aperture channel]\label{def:aperture}
A channel $c$ of $P$ is an \emph{aperture channel} when the data on $c$ are supplied by
the environment --- in the context-labelled semantics of \cite{probing}, by the
context $K$ in $K[P]$ rather than by $P$ --- and no presentation of the possible
supplies is given.
\end{definition}

\begin{construction}[$\AC$ at an aperture]\label{cons:ac-e}
\begin{lstlisting}
for( x <- sensor  where  psi(x) ) P
\end{lstlisting}
where \texttt{sensor} is an aperture channel. $\Cand$ is indexed by whatever the
environment supplies; the clause's support carves out a subfamily of it; and a resolver
must select from a family over an index set with no presentation. That is full $\AC$.
\end{construction}

This is not a trick. It is \cite{agency}'s own thesis rendered syntactically. That
essay's central structural claim is that a confluent interior carries no choice
resource, that ``the resource degree of a macro-component is carried entirely by its
boundary'', and that ``capacity, like choice, lives at the apertures''. Route E says
the same thing about clauses: \emph{the only place in a program where full $\AC$ can be
demanded is a boundary at which the world, not the term, supplies the contention.} A
shard's oracle-consuming sites are exactly its I/O sites, and they are syntactically
identifiable.

\begin{remark}[This is also a security statement]
If full $\AC$ is reachable only at apertures, then a static analysis that identifies
aperture channels bounds the choice type of an entire shard. The unbounded altitude of
a deployment is confined to its declared boundary, which is precisely the discipline a
capability system already wants to impose for other reasons.
\end{remark}

\subsection{Route R: reflection, and Krivine's \texttt{quote}}\label{ssec:reflection}

The second hypothesis of Proposition~\ref{prop:ceiling} is that $\psi$ contains no
$\langle @ \rangle_\V$ under a fixed point. Drop that instead, and the family becomes
impredicative.

\begin{construction}[$\AC$ by reflection]\label{cons:ac-r}
\[
  \psi \;=\; \nu X.\, \big( \langle @ \rangle_\V X \tensorv \crisp{\beta(*y)} \big)
\]
at a site binding $y$. The graded namespace modality evaluates $X$ at the process
$*y$ quoted by the received name. Since names are quotations of processes and processes
may carry clauses, the domain over which $\langle @ \rangle_\V$ ranges includes
processes whose own clauses are instances of $\psi$. The fixed point is therefore taken
over a domain whose description mentions the fixed point being taken.
\end{construction}

\begin{lstlisting}
for( y <- c  where  nu X. ( <@>_V X (x) [beta(*y)] ) ) P
\end{lstlisting}

The family here is not presented in advance of the run, and not presented \emph{by} the
run either: it is specified by a condition quantifying over a domain that contains the
condition. This is the shape of full $\AC$ as \cite{choicesched} describes it --- an
unindexed contention surveyed coherently --- rather than of any $C_X$ for a fixed
represented $X$.

\begin{conjecture}[Degree of the reflective clause]\label{conj:refl}
Construction~\ref{cons:ac-r} is not Weihrauch-reducible to $C_{\Bai}$; its resolver
requires a choice function on an arbitrary family, i.e.\ full $\AC$.
\end{conjecture}

We flag this as a conjecture rather than a proposition because making it precise
requires a representation of the impredicative family, and the natural candidates
either collapse it back to $C_{\Bai}$ by stratifying the namespace, or fail to be
representations at all. The stratification question is the same one that returns in
\S\ref{sec:sr}, which we take as evidence that it is the real content rather than a
technicality.

There is, however, a striking piece of external corroboration. Krivine's classical
realizability interprets dependent choice using a \texttt{quote} instruction --- a
primitive that assigns a code to a term, breaking the purely functional discipline of
the abstract machine \cite{krivine}. \cite{choicesched} already notes this as ``the
cleanest existing statement that a choice-type's inhabitant is a program with matching
extra power''. The observation to add here is that

\begin{quote}
\emph{the rho calculus has \texttt{quote} as a term former.} What Krivine must add to
his machine to realise choice, $@P$ already is.
\end{quote}

\noindent
Reflection is not an incidental convenience of the rho calculus that happens also to
enable Construction~\ref{cons:ac-r}. It is, on this reading, the calculus's native
choice primitive, and it has been sitting in the syntax since \cite{Mer05}.

\subsection{Exactly two routes}\label{ssec:tworoutes}

\begin{corollary}[Trichotomy]\label{cor:trichotomy}
Let $P$ be a term and $\psi$ a clause at a site of $P$. Exactly one of the following:
(i) $P$ is closed with presented channels and $\psi$ has no $\langle @ \rangle_\V$
under a fixed point, and then $\chi \lw C_{\Bai}$ by Proposition~\ref{prop:ceiling};
(ii) the site's channel is an aperture, and $\chi$ may be $\AC$ by
Construction~\ref{cons:ac-e}; (iii) $\psi$ is reflective, and $\chi$ may be $\AC$ by
Construction~\ref{cons:ac-r}.
\end{corollary}

The corollary is what makes the ceiling useful rather than discouraging. It is not that
$\AC$ is unreachable; it is that the two ways of reaching it are both syntactically
conspicuous. One is an I/O boundary and the other is a clause that quotes. A checker
that flags those two features flags every site in a program that can demand a genuine
choice function.

\subsection{The Diaconescu cost}\label{ssec:diaconescu}

Diaconescu's theorem says that in a topos the axiom of choice implies the law of
excluded middle. The analogue here is a cost, and it is worth stating because it is the
price of the answer to the question this section asks.

\begin{principle}[$\AC$ and interference do not cohabit]\label{prin:diaconescu}
Suppose a resolver realises full $\AC$ over the graded families a language of clauses
defines: every such family admits a section produced by that resolver. Then the
resolver's outcome at every site is a single candidate --- a point --- so the
coefficient it assigns to any outcome is the value of the clause at one candidate, and
no cancellation between candidates can occur. Hence the resolver exhibits no
interference deficit, and it is insensitive to any structure in $\V$ beyond the support
of the clause. An $\AC$-realising resolver reads its clauses crisply.
\end{principle}

\begin{proof}[Sketch, and its limits]
A choice function selects an element of the fibre. The interference phenomenon of
\S\ref{sec:transverse} is the assignment of coefficient $0$ to an outcome reachable
along two candidates with nonzero clause values; that requires the resolver to sum
across candidates rather than select one. The two are exclusive by definition of
``section''. What is \emph{not} established here is the converse Diaconescu direction
--- that the availability of choice forces the crisp sort's excluded middle as a
theorem of the logic --- which needs the subobject structure of a topos of graded
predicates that we have not built. We state the forward direction only.
\end{proof}

\begin{remark}[Why this is a feature]
The two research programmes this note joins want different things: \cite{choicesched}
wants the strongest choice principles typed and bounded, \cite{gradedwhere} wants
interference arranged deliberately. Principle~\ref{prin:diaconescu} says these are
demands on different resolvers, not competing demands on one. A shard may run an
$\AC$-strength resolver at its apertures and a coherent resolver in its interior; what
it may not do is expect one resolver to be both. The next section says why in the
lattice's own terms.
\end{remark}

%%%%%%%%%%%%%%%%%%%%%%%%%%%%%%%%%%%%%%%%%%%%%%%%%%%%%%%%%%%%%%%%%%%%%%%%%%%%%%%
\section{Interference is transverse to the lattice}\label{sec:transverse}

The complex row of the algebra table is not a high altitude. It is not an altitude at
all.

\begin{definition}[Interference deficit \cite{gradedwhere}]\label{def:deficit}
The \emph{interference deficit} of a term $P$ is
$\mathrm{Reach}(P) \setminus \mathrm{supp}(\sem{P})$: the outcomes operationally
reachable from $P$ that carry coefficient $0$ in the denotation.
\end{definition}

\begin{proposition}[A deficit is not a section]\label{prop:transverse}
If the interference deficit of $P$ at a site is non-empty, no selection function on
$\Cand$ produces the resolver's behaviour at that site.
\end{proposition}

\begin{proof}
A section of a family lands in a fibre: whatever a choice function returns is an
element of the set it was asked to choose from, and the outcome that element produces
therefore carries positive weight. An outcome in the deficit is reachable --- some
candidate produces it --- but carries coefficient $0$. So the resolver's behaviour is
not the image of any section.
\end{proof}

The mechanism is the one already worked out in \cite{gradedwhere}: two candidates whose
outcomes coincide contribute $w_0 \sumv w_1$ to that outcome, and over $\Cx$ the sum of
two nonzero terms can vanish. We recomputed the minimal case (\texttt{choicesim.py},
E5): with $|w_0| = |w_1| = 1/\sqrt2$ and relative phase $0$, $\pi$, $\pi/2$ the
coefficient is $1.4142$, $0.0000$, $0.7071{+}0.7071i$, giving $|{\cdot}|^2$ of
$2.000000$, $0.000000$, $1.000000$. Over $\Rnn$ this cannot happen: minimising
$w_0 + w_1$ over $100 \times 100$ strictly positive pairs gives $0.020000$, the
smallest pair, and never $0$.

\begin{principle}[Two axes, one crossing]\label{prin:axes}
The Weihrauch axis measures \emph{how hard it is to produce a witness}. The value-algebra
axis measures \emph{whether ``which candidate?'' has an answer}. Over $\Bool$, an
idempotent graded algebra, or $\Rnn$, it does, and the resolver is a section, possibly
randomised; the choice hierarchy applies and the schemas of \S\ref{sec:formulae}
classify it. Over $\Cx$ it does not: the resolver sums rather than selects, and
destructive interference \emph{deletes} a reachable witness, which no section can do.
The two axes cross at exactly one point, $\WWKL$, where the lift into distributions is
still a lift into sections-with-randomness.
\end{principle}

\begin{remark}[Consistency with Principle~\ref{prin:diaconescu}]
The two statements are the same fact from opposite ends. $\AC$ forces a section, which
forbids a deficit; a deficit forbids a section, which forbids $\AC$. The
top of the choice hierarchy and the interfering regime are disjoint, and neither is
above the other.
\end{remark}

%%%%%%%%%%%%%%%%%%%%%%%%%%%%%%%%%%%%%%%%%%%%%%%%%%%%%%%%%%%%%%%%%%%%%%%%%%%%%%%
\section{Fairness disciplines are clause combinators}\label{sec:fairness}

\cite{choicesched} places fairness on the typing side: a fairness discipline constrains
which schedulers inhabit a choice-type. In the graded setting it moves to the term
side, because a clause can starve a candidate all by itself.

\begin{definition}[Support fairness]\label{def:supportfair}
A run is \emph{weakly support-fair} when a candidate continuously in
$\mathrm{supp}(\sem{\psi})$ eventually fires, and \emph{strongly support-fair} when a
candidate infinitely often in the support eventually fires.
\end{definition}

\begin{proposition}[Zero is absorbing]\label{prop:absorbing}
A receipt whose clause evaluates to $0$ is outside the support, hence never fires,
hence never changes the term against which its clause is evaluated. If the clause's
value depends only on state that firing would update, the value remains $0$ forever.
\end{proposition}

This is not hypothetical. \cite{gradedwhere}'s forager exhibits it: with the PLN
confidence $c = n/(n+k)$ and no prior evidence, $n_0 = 0$ gives $c = 0$, gives clause
value $0$, gives a receipt that never fires and therefore never acquires the evidence
that would raise $n$. We reproduced the trap in a stripped-down model
(\texttt{choicesim.py}, E4): with $n_0 = 0$ and no repair, $0$ firings over the
horizon and a terminal stack of $0.00$; with $n_0 = 1$, $197$ firings and $89.20$.

The repair \cite{gradedwhere} found is a disjunct, and in the present language that is
literally a formula.

\begin{definition}[The exploration combinator]\label{def:explore}
For $\epsilon \in V$ with $\epsilon \neq 0$, write $\psi^{\epsilon} = \psi \sumv
\epsilon$.
\end{definition}

\begin{proposition}[Weak fairness is a combinator]\label{prop:weakfair}
Over an algebra with no additive inverses,
$\mathrm{supp}(\psi^\epsilon) = \Cand$, so every enabled candidate is in the support at every step and every run is weakly
support-fair.
\end{proposition}

\noindent
Measured: the same forager with $n_0 = 0$ and $\psi \sumv 0.02$ recovers to $179$
firings and a terminal stack of $83.80$, against $0$ and $0.00$ without it
(\texttt{choicesim.py}, E4). Exploration is a disjunct; fairness is a formula.

\begin{remark}[Strong fairness is a fixed-point condition, and it costs]
Weak support fairness is cheap because it is a condition on each step. Strong support
fairness is a condition on infinite runs --- ``infinitely often in the support implies
eventually fires'' --- which is a $\nu$ over the run tree, and by \S\ref{ssec:wkl} a
$\nu$ over an infinite finitely-branching tree is $\WKL$. So the classical folklore
that strong fairness is a substantially stronger assumption than weak fairness acquires
a degree: the gap is compactness.
\end{remark}

\begin{remark}[The additive-inverses caveat is the interference caveat]
Proposition~\ref{prop:weakfair} requires that adding $\epsilon$ cannot cancel. Over
$\Cx$ it can: $\psi \sumv \epsilon$ may vanish where $\psi$ did not. So the exploration
combinator is a fairness device in the $\Bool$, $\Rnn$ and PLN regimes and a
\emph{de-tuning} device in the coherent one. The same structure that buys interference
costs the cheap fairness repair.
\end{remark}

%%%%%%%%%%%%%%%%%%%%%%%%%%%%%%%%%%%%%%%%%%%%%%%%%%%%%%%%%%%%%%%%%%%%%%%%%%%%%%%
\section{Subject reduction, and reflection as the suspect}\label{sec:sr}

\cite{choicesched}'s central metatheorem is subject reduction, read as: \emph{scheduling
cannot bootstrap an oracle}. \cite{agency} makes the same question the hinge of its
factorisation --- if $\chi(c)$ can rise as a term runs, an interior cut can cross the
threshold and an island cracks. In the present setting the question becomes concrete
and, we suspect, has a negative answer for the full language.

\begin{definition}[Clause level]\label{def:level}
The \emph{level} $\ell(\psi)$ of a clause is its position in Table~\ref{tab:master},
read syntactically: the polarity and boundedness of its fixed points, the width of the
sites it ranges over, and whether it contains $\langle @ \rangle_\V$ under a fixed
point.
\end{definition}

\begin{question}[Subject reduction]\label{q:sr}
If $P \to P'$ by a firing at a site whose clause has level $\ell$, is every clause in
$P'$ of level $\leq \ell$?
\end{question}

\begin{proposition}[Reflection defeats it]\label{prop:srfails}
No. Let $Q$ be a process containing a receipt whose clause is
$\nu X.(\crisp{\beta} \tensorv \langle K \rangle_\V X)$, at level $\WKL$. Let
\[
  P \;=\; c!(Q) \;\mid\; \mathrm{for}(y \leftarrow c \ \mathrm{where}\
  \crisp{\mathrm{true}})\, {*}y .
\]
Every clause \emph{in $P$} is fixed-point-free, so $\ell(P) = \bot$. But
$P \to {*}@Q \scong Q$, and $\ell(Q) = \WKL$. The level rose across a reduction whose
own clause was trivial.
\end{proposition}

\begin{proof}
Immediate from $*@Q \scong Q$ and the definition of level as a syntactic property of
the clauses present.
\end{proof}

The counterexample is embarrassingly cheap, and that is the point: the mechanism is not
some subtle cascade but the basic move of the calculus. A name is a quoted process; a
process may carry clauses; dropping a received name installs them. \emph{Reflection
lets a low-altitude term receive a high-altitude one.} In \cite{agency}'s language, the
island cracks not because a boundary decision manufactured strength inside, but because
strength was mailed in.

\begin{remark}[The same suspect, three times]
Reflection is now the suspect in three separate open questions of this programme. It
is the conjectured source of length-asymmetric reconvergence, without which
interference in unmodified rho is purely constructive \cite{gradedwhere}; it is
Route R to full $\AC$ (\S\ref{ssec:reflection}); and it is the counterexample to
subject reduction here. The recurrence is not a coincidence: all three are consequences
of the namespace being able to carry arbitrary computational content, which is what
$@P$ is for. Krivine's \texttt{quote} is the same primitive and it too is what buys
choice and what breaks the standard metatheory.
\end{remark}

\subsection{The repair, and where $\AC$ ends up}

\begin{definition}[Stratified \texttt{where}]\label{def:strat}
A shard declares a level ceiling $L$. A clause \emph{elaborates at level $L$} when
(i) its own shape is at level $\leq L$ by Definition~\ref{def:level}, and (ii) every
occurrence of $*y$ in it, for $y$ bound by the site, is guarded by a level check on the
received process --- refusing at run time any received term carrying a clause above
$L$.
\end{definition}

\begin{proposition}[Subject reduction, stratified]\label{prop:srholds}
In the fragment where every clause elaborates at level $L$ and every drop is
level-checked, Question~\ref{q:sr} has a positive answer, and the choice type of a term
is a reduction invariant bounded by $L$.
\end{proposition}

\begin{proof}[Proof sketch]
The only way to introduce syntax not present in $P$ is substitution of a received name
followed by a drop; (ii) refuses drops that would raise the level; (i) bounds what is
present initially; and no other rule of the calculus creates clauses.
\end{proof}

\begin{corollary}[$\AC$ is outside]\label{cor:acoutside}
Route R (Construction~\ref{cons:ac-r}) is by definition unstratified, and Route E
(Construction~\ref{cons:ac-e}) violates the hypothesis of the ceiling proposition
rather than the stratification, so it survives stratification but only at declared
aperture channels. Hence in a stratified shard, full $\AC$ occurs exactly at declared
apertures, and the choice type of the interior is invariant.
\end{corollary}

\noindent
This is the answer we would want if we could choose it. The factorisation of
\cite{agency} is stable in the stratified fragment; $\AC$ is available, and available
precisely where that essay says capacity lives; and the failure mode is not hidden but
is a declared feature of a channel. Whether the stratified fragment is expressive
enough to be worth programming in is the obvious next question and we do not answer it.

%%%%%%%%%%%%%%%%%%%%%%%%%%%%%%%%%%%%%%%%%%%%%%%%%%%%%%%%%%%%%%%%%%%%%%%%%%%%%%%
\section{Implementation}\label{sec:impl}

What changes, concretely, in rholang 1.4 and MeTTaIL \cite{omnibus}.

\begin{enumerate}[label=I\arabic*.,leftmargin=2.6em,itemsep=3pt]
\item \textbf{The graded sort grows fixed points.} Definition~\ref{def:grammar}: add
$X$, $\mu X.\psi$, $\nu X.\psi$ and $\langle @ \rangle_\V \psi$ to the graded
condition grammar, with guardedness checked at elaboration. The crisp sort is
unchanged.
\item \textbf{Unfolding is budgeted by default.} Remark~\ref{rem:budget}: the existing
inference-token budget caps unfoldings, which keeps every clause at $\bot$ unless the
programmer explicitly buys altitude. This makes the safe default the cheap default.
\item \textbf{The resolver is an interface, not a fixed engine.} A resolver is a
$J$-algebra: given a clause's value vector, return a candidate (or a distribution over
candidates). The existing RSpace race, the Gillespie sampler of \cite{weightedgslt},
and a Viterbi or tropical optimiser are four instances. Attainment
(Definition~\ref{def:realise}) is the conformance law such an interface should carry,
in the style of the existing \texttt{oslf.rs} conformance-law idiom.
\item \textbf{Level as an elaboration-time judgement.} Definition~\ref{def:level} is
computed from the clause's syntax tree; Definition~\ref{def:strat} adds the drop check.
Both are decidable and neither requires running anything. This is the same shape as the
isometry judgement proposed in \cite{gradedwhere} for the coherent fragment, and the two
should share machinery.
\item \textbf{Aperture channels are declared.} Route E needs a channel attribute
distinguishing environment-fed channels from term-fed ones, which a capability
system plausibly wants anyway. Corollary~\ref{cor:acoutside} then makes the shard's maximal
choice type a static read-off.
\item \textbf{Candidate identity must be by occurrence.} As \cite{gradedwhere} already
notes, an implementation keying data by content, or allocating identifiers during
firing, cannot represent $\Cand$ faithfully --- and the width engine of
\S\ref{sec:engines} is precisely a statement about $|\Cand|$, so the addressing defect
would corrupt the level judgement as well as the values.
\end{enumerate}

%%%%%%%%%%%%%%%%%%%%%%%%%%%%%%%%%%%%%%%%%%%%%%%%%%%%%%%%%%%%%%%%%%%%%%%%%%%%%%%
\section{Open problems}\label{sec:open}

\begin{enumerate}[label=O\arabic*.,leftmargin=2.6em,itemsep=3pt]
\item \textbf{Degrees, properly.} Fix a representation of the candidate family as a
point of a represented space and prove that each schema of \S\ref{sec:formulae} has
the degree claimed, rather than merely imposing the corresponding obligation.
\item \textbf{Conjecture~\ref{conj:one}.} Prove that the Escard\'o--Oliva product, the
Weihrauch composition $\Wstar$, and chaining through $\langle K \rangle_\V$ agree, on
the attainable fragment. This is what makes $\DC$ the $\omega$-fold product and bar
recursion the scheduler.
\item \textbf{Conjecture~\ref{conj:refl}.} Settle the degree of the reflective clause.
Either exhibit a representation collapsing it to $C_{\Bai}$ --- in which case
Route R is not a route and $\AC$ is reachable only at apertures --- or show the
impredicativity is essential.
\item \textbf{Diaconescu, properly.} Build the topos of graded predicates over a GSLT
in which Principle~\ref{prin:diaconescu} is a theorem with its converse, rather than a
forward implication and an analogy.
\item \textbf{Expressiveness of the stratified fragment.} Proposition~\ref{prop:srholds}
buys subject reduction by refusing drops that raise the level. Is the resulting language
still expressive enough for the applications --- in particular, does it retain the
self-hosting and DSL-mobility properties that motivate reflection in the first place?
\item \textbf{Non-attainable resolvers as a monad.} \S\ref{ssec:lift} lifts the
selection function's codomain twice, to distributions and to amplitudes. Is there a
uniform account --- a parameterised selection monad over a choice of enrichment --- of
which the classical, randomised and coherent resolvers are the three instances?
\item \textbf{The altitude of consensus.} \cite{choicesched} proposes stacking its
pipeline with the Coalition-Logic consensus extraction. With the present schemas one
can ask a sharper question: what is the level, in the sense of
Definition~\ref{def:level}, of the clauses a correct-by-construction consensus protocol
requires, and is it below the ceiling a shard would want to declare?
\end{enumerate}

%%%%%%%%%%%%%%%%%%%%%%%%%%%%%%%%%%%%%%%%%%%%%%%%%%%%%%%%%%%%%%%%%%%%%%%%%%%%%%%
\section{Conclusion}\label{sec:conc}

The scheduler calculus that \cite{choicesched} lists as an open problem is, we propose,
the graded \texttt{where} clause of \cite{gradedwhere} with its fixed point restored. A
clause is a quantifier and a resolver is a selection function; the monad morphism
between them says when the resolver realises the clause; and attainability of the
induced quantifier is a sharp line, computable, and drawn exactly where the algebra's
addition stops being idempotent. Below that line the scheduler is a pure term. Above it
the scheduler must optimise, or randomise, or --- in the coherent regime --- stop being
a scheduler at all.

The ascent through the choice hierarchy is then the reintroduction of a single
syntactic feature, the graded fixed point, driven by three engines: depth, width, and
reflection. Depth and width are incomparable, which is why the order on scheduler types
is a lattice, and the incomparability is visible in the formula rather than in the
metatheory. Full $\AC$ is reachable, by exactly two routes, and both are conspicuous:
a channel fed by the world, or a clause that quotes. The first is where \cite{agency}
already said capacity lives. The second is what Krivine had to add to his machine and
what the rho calculus has had since 2005.

Two facts sit at the edges of the picture and both are, we think, load-bearing rather
than awkward. The complex-valued regime is transverse to the hierarchy, not high in it,
because a resolver that interferes deletes a reachable witness and no section can do
that. And the mechanism that buys full $\AC$ is the mechanism that breaks subject
reduction, which is \cite{agency}'s open hinge with a two-line counterexample and a
stratification that repairs it. Taken together they say that a shard can have any two
of unrestricted reflection, invariant choice types, and coherent resolution --- and
that the choice of which two is a language-design decision that can now be written
down, checked at elaboration, and read off a clause.

\begin{thebibliography}{99}

\bibitem[Mer05]{Mer05} L. G. Meredith and M. Radestock. A reflective higher-order
calculus. \emph{Electr.\ Notes Theor.\ Comput.\ Sci.}, 141(5):49--67, 2005.

\bibitem[Agcy]{agency} L. G. Meredith. \emph{Where Does Agency Live? Nondeterminism as
a computational resource, and agency across the computational and biological axes}.
Manuscript, F1R3FLY, 2026.

\bibitem[ChSch]{choicesched} L. G. Meredith. \emph{From the Axiom of Choice to Fair
Scheduling: Toward a Curry--Howard Correspondence for Nondeterminism, in which choice
principles are the types and fair schedulers are the terms}. Research Note, F1R3FLY,
2026.

\bibitem[GWhere]{gradedwhere} L. G. Meredith. \emph{Graded Where-Clauses: Non-crisp
truth values as the syntax of non-determinism resolution in the rho calculus, with
quantum and probabilistic-logic instances}. Research Note, F1R3FLY, 2026.

\bibitem[cwZX]{cwrhozx} L. G. Meredith. \emph{Complex-Weighted Rho to ZX: Compiling
weighted graph-structured lambda theories to the ZX-calculus by symbolic execution}.
Research Note, F1R3FLY, 2026.

\bibitem[WGSLT]{weightedgslt} L. G. Meredith. \emph{Weighted Graph-Structured Lambda
Theories: Stochastic and quantum execution for MeTTaIL}. Research Note, F1R3FLY, 2026.

\bibitem[Omni]{omnibus} L. G. Meredith. \emph{Graph-Structured Lambda Theories: A
Reference for the Working Developer}. F1R3FLY, 2026.

\bibitem[Probe]{probing} L. G. Meredith. \emph{Probing Universality: Context-Labelled
Bisimulation and a Relative Refinement of Turing Completeness for Interactive GSLTs}.
Research Note, F1R3FLY, 2026.

\bibitem[FRho]{fuzzyrho} L. G. Meredith. \emph{A quantale-graded rho calculus}.
Manuscript, F1R3FLY, 2026.

\bibitem[BGP]{weihrauch} Vasco Brattka, Guido Gherardi, and Arno Pauly. Weihrauch
complexity in computable analysis. In \emph{Handbook of Computability and Complexity in
Analysis}, Springer, 2021.

\bibitem[BGH]{brattkalasvegas} Vasco Brattka, Guido Gherardi, and Rupert H\"olzl.
Probabilistic computability and choice. \emph{Information and Computation},
242:249--286, 2015.

\bibitem[EO]{escardo} Mart\'in Escard\'o and Paulo Oliva. Selection functions, bar
recursion and backward induction. \emph{Mathematical Structures in Computer Science},
20(2):127--168, 2010.

\bibitem[Kri]{krivine} Jean-Louis Krivine. Dependent choice, `quote' and the clock.
\emph{Theoretical Computer Science}, 308(1--3):259--276, 2003.

\bibitem[Spec]{spector} Clifford Spector. Provably recursive functionals of analysis.
In \emph{Recursive Function Theory}, AMS, 1962.

\bibitem[ML]{martinlof} Per Martin-L\"of. \emph{Intuitionistic Type Theory}. Bibliopolis,
1984.

\bibitem[Dia]{diaconescu} Radu Diaconescu. Axiom of choice and complementation.
\emph{Proc.\ Amer.\ Math.\ Soc.}, 51:176--178, 1975.

\end{thebibliography}

\appendix
\section{The computations}\label{app:comp}

All numbers in this note are produced by \texttt{choicesim.py}, which has no
dependencies beyond the Python standard library. It runs five experiments.

\begin{description}[leftmargin=2.2em,itemsep=3pt]
\item[E1, attainability.] For each algebra and clause, tests over $20{,}000$ random
payoffs whether $\phi_\psi(p)$ is a value of $p$; then verifies
Proposition~\ref{prop:attain}(ii) by checking that ``attained'' agrees with ``the
winning candidate carries the unit'' ($100.00\%$); then verifies the distributional
lift by Monte Carlo over $400{,}000$ draws.
\item[E2, the product.] Computes the Escard\'o--Oliva product of two argmax selection
functions and the chaining of two graded sites through $\langle K \rangle_\V$ over
$20{,}000$ random payoff tables ($0$ disagreements), with a negative control that
chains in the wrong algebra ($4{,}483$ disagreements) so the agreement is not vacuous.
\item[E3, $\LLPO$.] Checks that the support of the two-sided race is never empty over
$3 \times 40 \times 60$ instance/stage pairs, and constructs, for each stage-$k$
selector with $k \leq 7$ and each tie-break, the diagonal instance defeating it.
\item[E4, fairness.] Runs the stripped-down forager in three configurations, exhibiting
the bootstrap trap ($0$ firings, terminal $0.00$) and its repair by $\psi \sumv
\epsilon$ ($179$ firings, terminal $83.80$) against the primed baseline ($197$,
$89.20$).
\item[E5, the deficit.] Computes the coefficient on the coinciding outcome at relative
phases $0$, $\pi$, $\pi/2$, and minimises the corresponding $\Rnn$ sum over
$100 \times 100$ strictly positive pairs to confirm it never vanishes.
\end{description}

\end{document}
